\documentclass[11pt]{article}
\usepackage[margin=1in]{geometry}
\usepackage{amsmath,amssymb,mathtools,booktabs}
\usepackage[hidelinks]{hyperref}
\newcommand{\diag}{\operatorname{diag}}
\newcommand{\tr}{\operatorname{tr}}
\newcommand{\TALO}{\mathrm{TALO}}
\newcommand{\ALO}{\mathrm{ALO}}
\newcommand{\E}{\mathbb E}
\newcommand{\R}{\mathbb R}

\title{\texttt{ppmltalo}: Estimator and Numerical Integration Notes}
\author{PPML TALO project}
\date{Development version, 13 August 2026}

\begin{document}
\maketitle

\begin{abstract}
This note is the mathematical companion to the Stata/Mata command
\texttt{ppmltalo}.  It fixes the point-estimator contract, derives the complete
PPML target Hessian and analytic deletion formulas, and translates every dense
expression into matrix--vector operations.  It also records the separate
independent-cluster formula, its exact frequency-collapsed representation, and
both randomized and deterministic coefficient-basis trace calculations.  The
software reports point estimates only.  Randomized probe error is numerical
error, not an econometric standard error.
\end{abstract}

\section{Model, quotient, and target}

Let $Y\in\R^n$ have conditional mean
\[
 \mu(\beta)=\exp(o+X\beta),
\]
where the exponential is componentwise.  Fixed effects make $X$ rank
deficient in a redundant parameterization.  All inverse and downdate formulas
below are evaluated on a fixed full-column-rank quotient of the
target-admissible face.  A Moore--Penrose inverse in a redundant
parameterization is only an invariant interpretation of the resulting
observation-space objects, under fixed rank, a common identified range, and
right-hand sides in that range.  Released code uses an explicit quotient and
never adds a ridge penalty.

At the PPML estimate $\widehat\beta$, define
\[
 \widehat\eta=o+X\widehat\beta,\qquad
 \widehat\mu=\exp(\widehat\eta),\qquad
 \widehat r=Y-\widehat\mu,\qquad X'\widehat r=0,
\]
\[
 W=\diag(\widehat\mu),\qquad H=X'WX,\qquad A=H^{-1},
 \qquad C=XAX'.
\]
The target is the identified quadratic form
\[
 \theta_Q=\beta_0'Q\beta_0,
 \qquad Q=Q',
\]
with plug-in $T_Q(Y)=\widehat\beta'Q\widehat\beta$.  Identification requires
$Qv=0$ for every null direction $v$ of the relevant design.  This condition
also makes the target invariant to the arbitrary fixed-effect normalization.

\section{Complete target curvature}

For a perturbation $u$ of the outcome vector, implicit differentiation of the
score $X'\{Y-\mu(\beta)\}=0$ gives
\[
 D\widehat\beta[u]=AX'u.
\]
Put
\[
 s_Q=XAQ\widehat\beta,
 \qquad d_Q=\widehat\mu\circ s_Q.
\]
Differentiating the plug-in a second time shows that its observation-space
\emph{half-Hessian} is
\begin{equation}
 G_Q=XAQAX'-C\diag(d_Q)C.                         \label{eq:G}
\end{equation}
Thus the second-order Taylor term is $u'G_Qu$; there is no extra factor of
$1/2$ in the correction.  The first term in \eqref{eq:G} is the coefficient
covariance curvature.  The second is the nonlinear coefficient--mean
curvature.  Dropping the second term produces a linear-model correction, not
TALO.

Both $s_Q$ and $G_Q$ are linear in $Q$.  If
$Q_T=Q_W+Q_F+2Q_{WF}$, every probe-level implementation must obey the same
identity.  Common probes across targets are therefore part of the numerical
contract.

\section{Independent-row analytic TALO}

For row $i$, define its fitted leverage
\[
 h_i=\widehat\mu_i x_i'Ax_i=(W^{1/2}XAX'W^{1/2})_{ii}.
\]
Deleting row $i$ changes the information to
$H_{-i}=H-\widehat\mu_i x_ix_i'$.  A Newton step from the full-sample PPML
solution toward the deleted solution, combined with the Sherman--Morrison
identity, gives the one-step ALO index
\begin{equation}
 \widetilde\eta_{i,-i}^{\ALO}
 =\widehat\eta_i-
 \frac{h_i}{1-h_i}\frac{\widehat r_i}{\widehat\mu_i}.       \label{eq:singleton-index}
\end{equation}
The default exponential prediction is
$\widetilde\mu_{i,-i}^{\ALO}=\exp(\widetilde\eta_{i,-i}^{\ALO})$.
For a numerical full-score residual $s=X'\widehat r\ne0$, the Newton
coefficient step is $H_{-i}^{-1}(s-x_i\widehat r_i)$ and therefore contains
the additional term $H_{-i}^{-1}s$.  The command reports and gates a
normalized full-score residual before using \eqref{eq:singleton-index}.
The implemented statistic is
\begin{equation}
 \widehat\theta_Q^{\TALO}
 =\widehat\beta'Q\widehat\beta-
 \sum_{i=1}^n G_{Q,ii}Y_i
       (Y_i-\widetilde\mu_{i,-i}^{\ALO}).                   \label{eq:talo}
\end{equation}
A separate linearized diagnostic would replace the last product by
$Y_i\widehat r_i/(1-h_i)$.  It is not exposed by the current command.  The
exponential and linearized versions are different finite-sample statistics and
must be labeled separately if the diagnostic is added later.

The release contract withholds \eqref{eq:talo} after a noninterior fit, a
target-relevant quotient failure, a deletion-induced face or rank failure, a
nonfinite analytic prediction, or leverage outside the declared safe range.
The current prototype checks the retained full fit, worker--firm connectivity,
every row for a worker--firm multigraph bridge, every positive-outcome deletion
for loss of a positive worker or firm margin, operational solver checks,
predictions, and a leverage safety diagnostic.  The bridge check is a
linear-time depth-first low-link calculation and treats parallel observation
edges separately.  These deletion gates do not yet certify every
deleted-sample PPML face outside the supported pure two-way model.  The
positive-graph certificate below closes that finite face gate for the current
scope.  Formula applicability remains unverified because downstream
statistical-theory and scale gates remain open.  The implementation never clips
leverage and never substitutes ridge output.

For the supported pure two-way worker--firm model, the code uses the following
sufficient deleted-face certificate.  A Poisson separation direction has row
index $\delta_{wf}=a_w+b_f$, satisfies $\delta_{wf}=0$ on every
positive-outcome row, and satisfies $\delta_{wf}\leq0$ on every zero-outcome
row, with at least one strict inequality.  If the bipartite multigraph formed
by the positive-outcome rows spans every worker and firm and is connected, its
equalities imply $a_w=c$ and $b_f=-c$ for a common $c$.  Hence
$\delta_{wf}=0$ on every retained row, so no nontrivial separation direction
exists.  If that positive-outcome multigraph has no bridge, deletion of any
positive row preserves this spanning connectivity; deletion of a zero row
does not change the positive graph.  Thus every one-row-deleted fit remains on
the supported interior face.  In parallel, the full observation multigraph
must have no bridge so every deletion preserves quotient rank.  These are
conservative sufficient conditions for the two-way fixed-effect model, not a
necessary characterization and not a certificate for additional regressors or
fixed-effect dimensions.

For completeness, this sign characterization follows directly from the
finite-dimensional Poisson log likelihood.  For a coefficient direction $d$
with row index $delta=Xd$ and a fixed finite offset $o$,
\[
 \ell(\beta+td)-\ell(\beta)
 =t\sum_iY_i\delta_i-
   \sum_i\mu_i(\beta)\{\exp(t\delta_i)-1\},\qquad t\geq0.
\]
Avoiding divergence to minus infinity requires $\delta_i=0$ when $Y_i>0$
and $\delta_i\leq0$ when $Y_i=0$.  A strict inequality on any zero row makes
the likelihood increase toward a boundary limit and rules out a finite
maximizer.  Conversely, on a full-rank finite quotient, absence of such a
nonzero row-index direction makes the log likelihood coercive: any unbounded
coefficient sequence has a normalized subsequence whose limiting direction
would otherwise satisfy these recession inequalities.  Strict concavity then
gives a unique finite maximizer.  The offset affects the positive levels
$\mu_i(\beta)$ but not the directional indices and therefore does not alter
the recession cone.

Connected positive support forces the worker direction coefficients to a
common $c$ and the firm coefficients to $-c$.  This is the usual global
normalization shift in redundant coordinates and has zero row index.  On the
implemented quotient, which omits one firm level, the normalization forces
$c=0$.  Thus the certificate excludes both a nonzero separating index and a
nonzero direction on the selected quotient.

The certificate is deliberately stronger than the exact condition.  If the
positive graph has several spanning components, its equalities leave one
constant per component.  Each zero row gives a directed inequality from its
worker component to its firm component.  Conditional on full observation-
graph connectivity, the exact finite-face condition is strong connectivity
of this directed component graph.  The production gate does not implement
that weaker necessary-and-sufficient test; it may withhold valid samples.
The independent test oracle implements the component condition on exhaustive
small designs and verifies that every design accepted by the production gate
is safe under every one-row deletion.

\section{Independent clusters}

Let $\{c\}$ partition the rows into mutually independent clusters.  Denote
cluster restrictions by $X_c,W_c,Y_c,\widehat r_c$, and define
\begin{equation}
 P_c=W_c^{1/2}X_cAX_c'W_c^{1/2}.                     \label{eq:Pc}
\end{equation}
Deleted-sample rank stability on the quotient is equivalent to
$\lambda_{\max}(P_c)<1$.  With
\[
 a_c=W_c^{-1/2}\widehat r_c,
 \qquad q_c=(I-P_c)^{-1}a_c,
\]
and the exact full-sample score condition $X'\widehat r=0$, Woodbury gives an
exact representation of the one-Newton-step change of the fitted indices on
the deleted cluster,
\begin{equation}
 \Delta\eta_c^{\ALO}
 =-W_c^{-1/2}P_c(I-P_c)^{-1}a_c
 =-W_c^{-1/2}(q_c-a_c).                              \label{eq:cluster-shift}
\end{equation}
Consequently
\[
 \widetilde\mu_c^{\ALO}
 =\exp(\widehat\eta_c+\Delta\eta_c^{\ALO}),
\]
while the linearized deleted residual is
\begin{equation}
 \widehat r_{c,-c}^{\ALO,\mathrm{lin}}
 =W_c^{1/2}(I-P_c)^{-1}W_c^{-1/2}\widehat r_c.        \label{eq:cluster-resid}
\end{equation}
The cluster point correction is
\begin{equation}
 \widehat\theta_{Q,\mathrm{cl}}^{\TALO}
 =\widehat\beta'Q\widehat\beta-
 \sum_cY_c'G_{Q,cc}(Y_c-\widetilde\mu_c^{\ALO}).     \label{eq:cluster-talo}
\end{equation}
For singleton clusters, \eqref{eq:cluster-shift}--\eqref{eq:cluster-talo}
reduce exactly to \eqref{eq:singleton-index}--\eqref{eq:talo}.  The cluster
formula changes the point estimate; clustering is not merely a standard-error
choice.

The repository's sparse-network TALO expansion is stated for singleton
observations.  Equations \eqref{eq:cluster-shift}--\eqref{eq:cluster-talo} are
therefore an implementation formula, not a claim that the singleton
asymptotic theorem automatically transfers to arbitrary clusters.  Production
cluster mode remains gated by review, exact-refit tests, and explicit resource
limits.

\subsection{Exact frequency collapse for match deletion}

Suppose stored row $j$ represents $f_j\in\{1,2,\ldots\}$ identical physical
rows with common $(Y_j,x_j,o_j)$, and let
\[
 F=\diag(f_1,\ldots,f_m),\qquad
 W^F=\diag(f\circ\widehat\mu),\qquad
 H^F=X'W^F X,\qquad A^F=(H^F)^{-1}.
\]
The weighted score is $X'F(Y-\widehat\mu)=0$. Literal expansion of each stored
row $f_j$ times gives exactly the same score and information matrix. The base
kernel shared by every pair of physical replicas of stored patterns $j$ and
$k$ is
\begin{equation}
 G_Q^F=XA^FQA^FX'-C^F\diag\{f\circ\widehat\mu\circ s_Q^F\}C^F,
 \quad C^F=XA^FX',\quad s_Q^F=XA^FQ\widehat\beta .       \label{eq:Gfreq}
\end{equation}
The outer frequency factors do not belong inside $G_Q^F$: they enter when the
physical rows in a deletion block are summed.
Thus $G_Q^F$ is the base kernel shared by pairs of physical replicas. If a
stored-coordinate perturbation changes all replicas of pattern $j$ together,
the half-Hessian in those stored coordinates is $FG_Q^FF$.

Let $c$ collect all stored patterns sharing one supplied deletion unit and
put $w_c^F=f_c\circ\widehat\mu_c$. The physical-row projection block and the
smaller collapsed block below have the same nonzero spectrum; the smaller
block may also contain zero modes,
\begin{equation}
 P_c^F=\diag(\sqrt{w_c^F})X_cA^FX_c'
       \diag(\sqrt{w_c^F}).                              \label{eq:Pfreq}
\end{equation}
Consequently the expanded deleted design keeps quotient rank if and only if
$\lambda_{\max}(P_c^F)<1$. Define
\[
 a_c^F=\sqrt{f_c/\widehat\mu_c}\circ
       (Y_c-\widehat\mu_c),\qquad
 q_c^F=(I-P_c^F)^{-1}a_c^F.
\]
The common fitted-index shift of every physical replica represented by stored
pattern $j\in c$ is
\begin{equation}
 \Delta\eta_c^{\ALO,F}
 =-\diag(\sqrt{w_c^F})^{-1}(q_c^F-a_c^F).               \label{eq:freq-shift}
\end{equation}
This follows either by applying Woodbury to
$H^F-X_c'\diag(w_c^F)X_c$ or by restricting the expanded block to its
duplicate-constant subspace. Thus collapse does not approximate the analytic
deletion.

Collapse must be stratified by deletion identifier. A stored pattern cannot
represent physical copies assigned to different all-copy deletion units. If
otherwise identical physical rows belong to different clusters, they remain
separate stored rows. This is an estimator assumption, not a computational
convenience.

With $\widetilde\mu_c^{\ALO,F}=\exp(\widehat\eta_c+
\Delta\eta_c^{\ALO,F})$, define the stored-row physical masses
\[
 y_c^F=f_c\circ Y_c,\qquad
 e_c^F=f_c\circ(Y_c-\widetilde\mu_c^{\ALO,F}).
\]
The exactly collapsed point correction is
\begin{equation}
 \widehat\theta_{Q,\mathrm{match}}^{\TALO,F}
 =\widehat\beta'Q\widehat\beta-
   \sum_c (y_c^F)'G_{Q,cc}^F e_c^F .                    \label{eq:freq-correction}
\end{equation}
Every term in \eqref{eq:freq-correction} equals the corresponding sum over
literal physical replicas. In particular, one frequency factor appears on
each side of $G_{Q,cc}^F$.

Target mass is distinct from likelihood frequency. Let $\pi_j$ denote total
target mass represented by stored row $j$. The default is $\pi_j=f_j$; an
explicit \texttt{targetweight()} supplies $\pi_j$ directly and is not
multiplied by $f_j$. Distributing $\pi_j/f_j$ to each physical replica gives
the same centered target matrix $Q$ as the stored-row calculation.
The full-sample target masses, their normalization and centering, and hence
each $Q$ are frozen before any deletion. These formulas do not cover a target
that is renormalized or recentered on each deleted sample.

\subsection{One-copy observation deletion after frequency collapse}

Frequency collapse does not force whole-pattern deletion. Under
\texttt{deletion(observation)}, a stored pattern with frequency $f_j$ denotes
$f_j$ distinct but identical physical observations, and one leave-out action
removes exactly one physical copy. The full-sample information remains $H^F$,
but the downdate for one copy is
\[
 H^F_{-j,1}=H^F-\widehat\mu_jx_jx_j'.
\]
Consequently the relevant physical-copy leverage is
\begin{equation}
 h_j^{(1)}=\widehat\mu_jx_j'A^Fx_j,                    \label{eq:one-copy-h}
\end{equation}
not $f_jh_j^{(1)}$. At an exact weighted-score root, each identical copy has
the same analytic deleted prediction
\begin{equation}
 \widetilde\eta_{j,-1}^{\ALO,F}
 =\widehat\eta_j-
 \frac{h_j^{(1)}}{1-h_j^{(1)}}
 \frac{Y_j-\widehat\mu_j}{\widehat\mu_j},
 \qquad
 \widetilde\mu_{j,-1}^{\ALO,F}
 =\exp(\widetilde\eta_{j,-1}^{\ALO,F}).               \label{eq:one-copy-shift}
\end{equation}
The literal physical-row correction therefore collapses to
\begin{equation}
 \widehat\theta_{Q,\mathrm{obs}}^{\TALO,F}
 =\widehat\beta'Q\widehat\beta-
 \sum_j f_jG_{Q,jj}^F Y_j
       (Y_j-\widetilde\mu_{j,-1}^{\ALO,F}).            \label{eq:one-copy-talo}
\end{equation}
There is one outer factor $f_j$ in \eqref{eq:one-copy-talo}. Using the
collapsed-pattern leverage $f_jh_j^{(1)}$ would instead delete every copy;
using $f_j^2$ in the diagonal correction would likewise compute the
whole-pattern block contraction rather than leave-one-physical-observation
out. Thus observation and cluster deletion deliberately differ whenever any
$f_j>1$.

This distinction also changes the matrix-free leverage sketch. Let
$W=\diag(f\circ\widehat\mu)$ and
$P_W=W^{1/2}X(H^F)^{-1}X'W^{1/2}$. Then $P_W$ is an orthogonal projection and
\[
 h_j^{(1)}=\frac{(P_W)_{jj}}{f_j},\qquad
 \E_z\left\{\frac{(P_Wz)_j^2}{f_j}\right\}=h_j^{(1)}.
\]
Thus the implementation sketches with $W^{1/2}$ on both sides and divides
the squared projected row by $f_j$. Replacing $W^{1/2}$ by
$\diag(\sqrt{\widehat\mu})$ before squaring would target the diagonal of a
squared nonprojection and is not a one-copy leverage estimator.

\subsection{Generic local block compression and rich-design gates}

Let $E_c$ select every global quotient coordinate that is active in deletion
block $c$, and write $B_c=X_cE_c$. This definition does not require the worker
or firm coordinates to be constant in the block. Put
\[
 K_c=E_c'A^FE_c,
 \qquad R_c=B_c'\diag(w_c^F)B_c.
\]
Because $A^F$ is positive definite on the selected quotient, $K_c$ is a
positive-definite principal submatrix. The exact block projection is
\begin{equation}
 P_c^F=\diag(\sqrt{w_c^F})B_cK_cB_c'
       \diag(\sqrt{w_c^F}).                            \label{eq:generic-local-p}
\end{equation}
The nonzero eigenvalues of \eqref{eq:generic-local-p} equal those of
$K_c^{1/2}R_cK_c^{1/2}$. Moreover, if
$g_c=B_c'\{f_c\circ(Y_c-\widehat\mu_c)\}$, the deleted index shift is
\begin{equation}
 \Delta\eta_c^{\ALO,F}
 =-B_c(K_c^{-1}-R_c)^{-1}g_c.                          \label{eq:generic-local-shift}
\end{equation}
This is the coefficient-local Woodbury form used by the bounded engine.

An equivalent rank-revealing row-space construction sets
\[
 Z_c=\diag(\sqrt{w_c^F})B_c=U_cL_c,
\]
where $U_c$ is a thin orthonormal basis after removing exact local
dependencies. Since the range of $P_c^F$ lies in the range of $U_c$, define
$M_c=U_c'P_c^FU_c$. Then
\begin{align}
 \lambda_{\max}(P_c^F)&=\lambda_{\max}(M_c),             \label{eq:local-spectrum}\\
 q_c^F-a_c^F
 &=U_c(I-M_c)^{-1}M_cU_c'a_c^F.                          \label{eq:local-inverse}
\end{align}
Equations \eqref{eq:local-spectrum}--\eqref{eq:local-inverse} show how a future
rank-revealing implementation can replace a stored-pattern square inverse by
one whose dimension is the local design rank. The current bounded engine uses
the algebraically equivalent coefficient-local form
\eqref{eq:generic-local-shift}. It applies its work limit to the number of raw
active quotient coordinates before computing and reporting the weighted local
rank. This is conservative: a block with many locally dependent active
coordinates can be withheld even though a thinner exact basis exists. It does
not truncate a small identified eigenvalue of $I-M_c$.
For the convenient legacy match case, constant worker and firm coordinates
make $B_c$ especially small. The generic construction gives the same answer
while also admitting supplied clusters that contain several worker or firm
coordinates.

For a possible randomized construction, define
\[
 P^F=(W^F)^{1/2}XA^FX'(W^F)^{1/2},
\]
and draw a Rademacher vector $z$. With
$u=P^Fz$ and $u_c$ denoting the restriction of this \emph{full projected
probe} to block $c$,
\[
 \E\{(U_c'u_c)(U_c'u_c)'\mid X,W^F\}=M_c.
\]
Using the block-local product $P_c^Fz_c$ would instead have expectation
$(P_c^F)^2$ and is incorrect. The implemented bounded design does not use
this randomized block construction: it forms exact coefficient-space
information inverses and exact low-dimensional block spectra. Randomization
is confined to the final trace contraction.
Nested probe counts and independent seeds measure only the final randomized
trace contraction. The deletion means, full block spectra, and positive block
spectra are fixed exact dense calculations in the bounded engine. The
configured full-block diagnostic must remain below the block limit and is not
an econometric confidence bound.

Graph connectivity alone no longer certifies rank after nuisance terms are
added. The full-design block spectrum provides the numerical rank gate. A
conservative interior-face certificate uses the positive-outcome information
\[
 W_+=\diag(f\circ Y),\qquad H_+=X'W_+X,\qquad A_+=H_+^{-1}.
\]
For block $c$, define
\[
 P_{+,c}=\diag(\sqrt{f_c\circ Y_c})X_cA_+X_c'
           \diag(\sqrt{f_c\circ Y_c}),
 \qquad
 H_{+,-c}=H_+-X_c'\diag(f_c\circ Y_c)X_c .
\]
Zero-outcome rows give zero rows and columns in $P_{+,c}$. If $H_+$ has full
quotient rank and every exact $P_{+,c}$ has spectral radius below one, then
$H_{+,-c}$ is positive definite and the positive rows after every deletion
span the full quotient. Any recession direction that is zero on those
positive rows therefore has zero row index everywhere, excluding a nontrivial
Poisson separation direction. A compressed implementation would need a
separate positive-weight local basis when zero weights change the local
column space. The bounded engine instead computes the exact spectrum from
$A_+$ and the separately weighted local coordinate map, so it does not reuse a fitted-mean
basis. This sufficient certificate may reject finite deleted fits that exist;
it must not be silently weakened. Exact small-design rank checks and deleted
PPML refits remain independent oracles.

All dense rank and inverse calculations use diagonal equilibration. For a
symmetric information matrix $H$ with positive diagonal, put
\[
 D_H=\diag\{\sqrt{H_{11}},\ldots,\sqrt{H_{pp}}\},\qquad
 S_H=D_H^{-1}HD_H^{-1}.
\]
The code tests the rank of $S_H$, computes $S_H^{-1}$, checks
$\lVert S_HS_H^{-1}-I\rVert_F/\sqrt p$, and returns
$H^{-1}=D_H^{-1}S_H^{-1}D_H^{-1}$. The same construction is applied to
$H_+$ and to the local kernel and deletion systems. Because every $D_H$ is
invertible, this changes coordinates only: rank, the block spectra, the
deleted index shift, and each target value are unchanged. It prevents a
large control-unit disparity from making a full-rank quotient appear singular
under a raw absolute scale. A nonpositive diagonal, deficient equilibrated
rank, nonfinite inverse, or failed equilibrated residual still produces a
typed withholding status.

\section{Matrix-free evaluation}

The production engine exposes four operations:
\[
 \mathcal X(v)=Xv,\qquad
 \mathcal X'(z)=X'z,\qquad
 \mathcal S(b)=H^{-1}b,\qquad
 \mathcal Q(v)=Qv.
\]
For fixed effects, $\mathcal X$ is a sum of group-level vectors and
$\mathcal X'$ is a set of grouped sums.  The information action used by
preconditioned conjugate gradients is
\[
 v\longmapsto \mathcal X'\{\widehat\mu\circ\mathcal X(v)\}.
\]
No observation-by-observation matrix is formed.

The worker block is especially large in an AKM design, but its columns are
mutually exclusive indicators.  Partition $X=(X_1,X_r)$, where $X_1$ is the
full worker block, and write
\[
 H=\begin{pmatrix}D&B\\B'&R\end{pmatrix},\qquad
 D=X_1'WX_1,\quad B=X_1'WX_r,\quad R=X_r'WX_r.
\]
Here $D$ is positive diagonal.  The production solver applies PCG to the
Schur complement
\[
 S=R-B'D^{-1}B,
 \qquad
 Sv=X_r'W\{I-X_1D^{-1}X_1'W\}X_rv.
\]
If $H$ is symmetric positive definite on the selected quotient, then $S$ is
symmetric positive definite.  In the pure two-way model the implementation
computes its exact diagonal from grouped worker--firm cell weights:
\[
 S_{ff}=\sum_{i:f(i)=f}\widehat\mu_i-
 \sum_w
 \frac{\left(\sum_{i:w(i)=w,\,f(i)=f}\widehat\mu_i\right)^2}
 {\sum_{i:w(i)=w}\widehat\mu_i}.
\]
It uses this positive diagonal as the Jacobi preconditioner.
For a right-hand side $b=(b_1',b_r')'$, it solves
\[
 Sx_r=b_r-B'D^{-1}b_1,\qquad
 x_1=D^{-1}(b_1-Bx_r).
\]
Every action remains a sequence of row maps and grouped sums. Batched columns
share those grouped operations. For nonzero $b$, exact reconstruction gives
\[
 b-Hx=\begin{pmatrix}0\\b_r-B'D^{-1}b_1-Sx_r\end{pmatrix}.
\]
The reduced stopping rule is therefore scaled by the norm of the original
$b$. An exactly zero right-hand side returns the zero solution by convention.
The implementation normalizes each right-hand side by its largest absolute
entry, uses scaled sums of squares and quad-precision inner-product
accumulation, and applies a stricter internal stopping margin. Every 100
iterations it replaces the recursive residual by an explicit Schur residual
without restarting the independent PCG recurrences. Batched columns retain
separate statuses and iteration counts, so a breakdown in one column does not
terminate healthy columns; an aggregate solve succeeds only when every column
succeeds. Acceptance uses a freshly recomputed double-precision full-system
residual $\lVert b-Hx\rVert_2/\lVert b\rVert_2$ on the normalized system.
This is an operational floating-point check. It is not a rigorous upper bound
on the exact-real residual, and a small residual does not guarantee small
forward error in an ill-conditioned system.

For any column vector $z\in\R^n$, define
\[
 u=\mathcal C(z)=\mathcal X\mathcal S\mathcal X'(z).
\]
Equation \eqref{eq:G} gives the exact action
\begin{equation}
 \mathcal G_Q(z)
 =\mathcal X\mathcal S\mathcal Q\mathcal S\mathcal X'(z)
  -\mathcal C\{d_Q\circ\mathcal C(z)\}.               \label{eq:Gaction}
\end{equation}
Each call uses only grouped reductions, target actions, and information solves
that pass the operational residual check.

\subsection{Exact low-dimensional nuisance augmentation}

Write the design as $X=(F,Z)$, where $F$ contains the two high-dimensional
target fixed-effect blocks and $Z$ contains a bounded number of jointly
estimated controls and nuisance fixed-effect quotient coordinates. Partition
the information as
\[
 H=\begin{pmatrix}H_F&H_{FZ}\\H_{ZF}&H_Z\end{pmatrix}.
\]
The matrix-free two-FE solver supplies actions of $H_F^{-1}$. The augmentation
uses a small number of those actions to construct
\begin{equation}
 V=H_F^{-1}H_{FZ},\qquad
 S_Z=H_Z-H_{ZF}V.                                      \label{eq:nuisance-schur}
\end{equation}
For a right-hand side $b=(b_F',b_Z')'$, define
\begin{align}
 t&=H_F^{-1}b_F,\\
 x_Z&=S_Z^{-1}(b_Z-H_{ZF}t),\\
 x_F&=t-Vx_Z.                                          \label{eq:nuisance-action}
\end{align}
Then $(x_F',x_Z')'=H^{-1}b$ exactly. The implementation diagonally
equilibrates $S_Z$, checks its rank and inverse residual, and then recomputes
the complete residual $b-Hx$ through implicit design actions. The same
construction is built independently for
\[
 H_+=X'\diag(f\circ Y)X,
\]
so the positive-face gate includes nuisance directions rather than checking
only the two target fixed effects. Here ``independently'' means a new weighted
information calculation on the \emph{same full-sample coefficient quotient}.
All stored rows and all full-sample nuisance levels remain in the design, with
zero positive weight when $Y_i=0$. The implementation never re-encodes
nuisance levels or chooses a new anchor after selecting positive rows. Such
re-quotienting can silently discard a positive-information null direction.
The implementation solves the columns of $V$ at a tolerance no looser than
$10^{-12}$ and forms both $H_{ZF}V$ and $V'H_FV$. These expressions are
identical when the base solves are exact. Their scaled discrepancy must be at
most $\max(10^{-10},100\epsilon_{\rm prep})$; otherwise the command returns a
typed Schur-coherence failure. This is an operational double-precision check,
not an exact-real forward-error theorem.
The conditional nuisance information is also compared with its unresidualized
weighted scale; a conditional share at or below the registered rank tolerance
is classified as nuisance collinearity before inversion.

Every occurrence of $A$ in the target half-Hessian, leverage, deleted index,
and trace action uses \eqref{eq:nuisance-action}. Although $Q$ has zero rows
and columns on nuisance coordinates, $AQ\widehat\beta$ generally has nonzero
nuisance coordinates because $A$ is not block diagonal. Hence setting the
nuisance coefficients' target action to zero does not justify omitting their
covariance from the inverse-information action.

For observation deletion the joint leverage splits exactly as
\[
 d_i x_i'H^{-1}x_i
 =d_i F_iH_F^{-1}F_i'
  +d_i\{Z_i-F_iV\}S_Z^{-1}\{Z_i-F_iV\}',
\]
where $d_i=\widehat\mu_i$ for one-copy deletion. The second term is computed
from the small residualized nuisance block. The first is computed exactly when
the two-FE quotient is bounded; otherwise effective resistance in any weighted
spanning tree supplies a deterministic upper bound because deleting off-tree
conductances can only increase resistance. Their sum is required to lie below
the unchanged leverage limit before the randomized pointwise sketch is used.
Before applying that gate, the implementation adds a nonnegative numerical
guard equal to the maximum of $10^{-10}$ and one hundred times the recorded
Schur-preparation and small-inverse residuals. Thus residual allowance can
only make the gate more conservative. The tree inequality is deterministic;
the complete floating-point certificate remains an operational numerical
certificate rather than a proved exact-real upper bound.
The sketch-plus-rowwise-MCSE diagnostic remains an additional numerical check,
not the deterministic certificate or a simultaneous probability bound.

The optional fixed-nuisance diagnostic first estimates the complete PPML
model, retains $Z\widehat\gamma$ in the fitted offset, and applies correction
actions only through $H_F^{-1}$. It never residualizes $Y$ and does not allow
$\widehat\gamma$ to move under deletion. This generally differs from
\eqref{eq:nuisance-action}. Even $H_{FZ}=0$ in the full-sample information is
insufficient for equality: deletion can create a nuisance score, move nuisance
coefficients, and change the deleted fitted index. Equality is action- and
deletion-specific and requires the complete joint and conditional row-space
actions to coincide for every curvature and deletion right-hand side. The
command therefore labels it
\[
 \texttt{CONDITIONAL\_ON\_ESTIMATED\_NUISANCE\_INDEX},
\]
not complete joint TALO.

\subsection{Randomized trace contraction}

Let $z$ have independent Rademacher entries, so $\E(zz')=I$.  For any square
matrix $M$,
\[
 \E(z'Mz)=\tr(M).
\]
For independent rows put
$K=\diag\{Y_i(Y_i-\widetilde\mu_i^{\ALO})\}$.  Then the correction in
\eqref{eq:talo} is
\begin{equation}
 \tr(G_QK)=\E_z\{z'\mathcal G_Q(Kz)\mid Y,X,K\},       \label{eq:row-trace}
\end{equation}
where the current trace probe is independent of the fitted data and of every
random sketch used to construct $K$.
The implementation does not apply $\mathcal G_Q$ separately for every target.
For each probe it solves
\[
 u=A X'z,\qquad v=A X'Kz.
\]
Since $A$ and $Q$ are symmetric, the same scalar is
\begin{equation}
 z'G_QKz
 =u'Qv-\sum_i d_{Q,i}(Xu)_i(Xv)_i.                    \label{eq:shared-probe}
\end{equation}
The two information solves and the two row actions in
\eqref{eq:shared-probe} are target independent.  Only the inexpensive $Qv$
action and the dot product with $d_Q$ change across targets.  Worker, firm,
covariance, and total corrections therefore use a common probe contribution,
which preserves the accounting identity before averaging.

For clusters let $e_c=Y_c-\widetilde\mu_c^{\ALO}$ and define the symmetric
block matrix
\begin{equation}
 K_c=\tfrac12(e_cY_c'+Y_ce_c'),\qquad K=\operatorname{blockdiag}(K_c). \label{eq:Kcluster}
\end{equation}
Since $G_Q$ is symmetric,
\[
 \tr(G_QK)=\sum_cY_c'G_{Q,cc}e_c.
\]
The cluster product needed by a probe is local:
\begin{equation}
 K_cz_c=\tfrac12\{e_c(Y_c'z_c)+Y_c(e_c'z_c)\}.          \label{eq:Kaction}
\end{equation}
Thus \eqref{eq:row-trace} remains valid with \eqref{eq:Kaction}.  Online
Welford accumulation reports the probe mean and its numerical Monte Carlo
standard error.  It is not an econometric standard error.

\subsection{Deterministic coefficient-basis trace}

The bounded match engine already stores the dense coefficient-space inverse
$A$, while its row count can remain in the millions.  It can therefore remove
the randomized trace error without constructing $X$, $K$, or $G_Q$ as dense
observation-space matrices.  Put $B=X'KX$, let $e_j$ be coefficient-space
standard basis vector $j$, and define
\[
 u_j=Ae_j,\qquad v_j=ABe_j.
\]
Cyclic invariance of the trace and \eqref{eq:G} give the exact identity
\begin{equation}
 \tr(G_QK)
 =\sum_{j=1}^p\left\{
   u_j'Qv_j-
   \sum_{i=1}^n d_{Q,i}(Xu_j)_i(Xv_j)_i
 \right\}.                                             \label{eq:basis-trace}
\end{equation}
Indeed, the first sum is
$\tr(AQAB)=\tr(XAQAX'K)$.  The second is
$\tr\{AX'\diag(d_Q)XAB\}=\tr\{C\diag(d_Q)CK\}$.

The code evaluates \eqref{eq:basis-trace} in batches.  For a batch $E$ of
standard basis columns it obtains
\[
 XE,\qquad BE=X'K(XE),\qquad U=AE,\qquad V=ABE,
\]
using only implicit design and block-kernel actions.  It then applies each
target $Q$ to $V$ and accumulates the two terms in
\eqref{eq:basis-trace}.  It stores the $p\times p$ information inverse already
required by match mode and only $O\{(n+p)b\}$ additional batch storage for
effective batch width $b=\min\{\texttt{batch},p\}$.  It never stores an
$n\times n$ matrix.  A configured batch at least as large as $p$ does create
the complete $n\times p$ direction block, so large-data runs should keep
\texttt{batch} small.  The calculation uses exactly $p$ coefficient
directions, is invariant
to the random seed, and reports zero trace Monte Carlo error.  Its work grows
linearly with $p$ matrix-free trace directions in addition to the dense
coefficient-space algebra, so the randomized trace remains available when a
declared numerical tolerance can be met more cheaply.  Inner products use
quad-precision accumulation and the streamed direction totals use compensated
summation.  The worker, firm, and covariance totals are combined explicitly
after accumulation to preserve the accounting identity.  Zero trace Monte
Carlo error means that no randomized trace contraction remains; it does not
bound floating-point or econometric error.

\subsection{Leverage and cluster-block sketches}

Let
\[
 P=W^{1/2}XAX'W^{1/2}.
\]
With exact application of $H^{-1}$ on the identified quotient,
$P=P'=P^2$.  For a Rademacher vector $z$, compute $u=Pz$ with one exact
information solve.  Then
\begin{equation}
 \E_z(u_iu_j\mid X,W)=P_{ij}.                            \label{eq:Psketch}
\end{equation}
In particular, $R^{-1}\sum_{r=1}^Ru_{ir}^2$ is a nonnegative estimate of
$h_i=P_{ii}$, and
$R^{-1}\sum_ru_{c,r}u_{c,r}'$ estimates the entire within-cluster block
$P_c$.

With an approximate solve producing an operator $\widetilde P$, the exact
second moment is $\widetilde P\widetilde P'$, not $P$.  Operational solver
residual checks control an approximation in the implemented floating
arithmetic; they do not make
\eqref{eq:Psketch} exact.

The nonlinear maps $(1-h_i)^{-1}$ and $(I-P_c)^{-1}$ make plug-in sketch error
nonlinear.  Unbiasedness of \eqref{eq:Psketch} alone does not make the final
TALO statistic unbiased in the numerical randomization.  Released code must
use independent sketch and trace streams, report convergence across nested
  probe counts, and withhold when a declared numerical safety diagnostic cannot
  keep the leverage or block spectrum away from one.  The current row-level
  implementation uses the estimate plus $3.290527$ rowwise Monte Carlo standard
  errors as an engineering gate; this quantity is not a calibrated simultaneous
  confidence bound.  Exact dense blocks remain the small-sample oracle.

Storing every $P_c$ costs $O(\sum_c|c|^2)$ doubles.  Cluster production mode
must declare a maximum cluster size and a total block-memory budget.  There is
currently no claim of $O(n)$ storage for unbounded cluster sizes.

\section{Target actions for an AKM decomposition}

Let $\pi_i\ge0$, $\sum_i\pi_i=1$, be target-population row weights, distinct
from the PPML likelihood weights.  For worker effects $\alpha_{w(i)}$ and firm
effects $\psi_{f(i)}$, define weighted means $\bar\alpha$ and $\bar\psi$.  The
four initial targets are
\begin{align*}
 \theta_W&=\sum_i\pi_i(\alpha_{w(i)}-\bar\alpha)^2,\\
 \theta_F&=\sum_i\pi_i(\psi_{f(i)}-\bar\psi)^2,\\
 \theta_{WF}&=\sum_i\pi_i(\alpha_{w(i)}-\bar\alpha)
                         (\psi_{f(i)}-\bar\psi),\\
 \theta_T&=\theta_W+\theta_F+2\theta_{WF}.
\end{align*}
If $Z_W$ and $Z_F$ map quotient coefficients to the row-level worker and firm
effects and $R_\pi=\diag(\pi)-\pi\pi'$, the corresponding symmetric coefficient
matrices are
\[
 Q_W=Z_W'R_\pi Z_W,\qquad Q_F=Z_F'R_\pi Z_F,\qquad
 Q_{WF}=\tfrac12\{Z_W'R_\pi Z_F+Z_F'R_\pi Z_W\}.
\]
Hence $\beta'Q_{WF}\beta$ is one covariance copy and
$Q_T=Q_W+Q_F+2Q_{WF}$.  The factor $1/2$ in $Q_{WF}$ and the factor $2$ in
$Q_T$ are both required.
Each $Q$ action is implemented by weighted group sums and centering.  The total
target is assembled by applying the displayed linear identity to every
plug-in, correction, and probe contribution. Nuisance fixed effects and
controls enter $H$ but receive zeros from these target actions. Observation
mode admits a bounded number of nuisance fixed-effect dimensions and numeric
controls through the common joint-information quotient. The dense cluster and
match engine applies the same full- and positive-information rank and face
gates to those nuisance coordinates.

\section{Numerical safeguards and returned evidence}

Every information solve returns its recomputed double-precision relative
residual, iteration count, and a typed termination status. Requested solver
tolerances below $10^{-15}$ are rejected, and 500 consecutive iterations
without a new best recursive residual produce a typed stagnation status. This is not a
rigorous exact-real residual or forward-error certificate. The release gate
requires all solves to converge,
all means and predictions to be finite and positive, the target to be
identified, the deletion unit to preserve the required quotient/face, and
leverage to remain below the configured limit with the required numerical
margin.  The prototype implements the stated conservative sufficient
deleted-face certificate for row deletion in a pure two-way model. The
frequency-weighted match extension uses the full-design and positive-design
block spectral gates above. In the bounded match engine these gates are exact
up to dense double-precision coefficient-space linear algebra; the configured
0.8 block limit is an additional conditioning threshold, not a statistical or
Monte Carlo bound. The code does not claim a weaker necessary-and-sufficient
rich-design face characterization. It records probe seed, count, batch size,
estimated numerical error, exact block diagnostics, equilibrated inverse
residuals, and reciprocal-condition estimates for dense information, local
systems, and the low-dimensional nuisance Schur block. A full condition number
is not formed for the large matrix-free information operator;
the benchmark harness separately records wall time.

When a whole-match gate rejects a calculation, the implementation retains the
ordinal block number and an aggregate stage code without exporting worker or
firm identifiers. The stages distinguish a nonpositive local
positive-kernel diagonal, deficient local positive-kernel rank, a
positive-deletion eigenvalue at one, a full-deletion eigenvalue at one, and
the stricter configured full-block threshold. Associated minimum and maximum
local metrics make a failed target-data run diagnostically useful while all
point-estimate matrices remain cleared.

The dense oracle, offset-aware exact deleted refits, finite-difference target Hessian,
matrix-free actions, exhaustive Rademacher identities on tiny samples, and
invariance tests are separate checks.  Their agreement verifies finite
algebra; it does not prove consistency or theorem applicability for an
empirical network.

\section{Status}

The finite formulas passed staged algebra, Schur, deleted-face, frequency,
block-Woodbury, and deterministic-trace gates. Each hard round received two
isolated GPT Pro reviews with frozen packets, verbatim records, and local
adjudication. The ST11 pairs reviewed independent deletion, joint nuisances,
and the implemented Schur/leverage/trace routing. Their adverse findings led
to the corrected weighted sketch, deterministic matrix-free release gates, a
common positive quotient, equilibrated inversion, final-finiteness checks,
forced accounting, post-fit cleanup without \texttt{e(V)}, collision-free pair
grouping, and condition diagnostics. One response contains contradictory
verdict blocks; both are preserved and the adjudication applies the adverse
disposition. The coefficient-local block engine reports local rank but still
applies its conservative raw active-coordinate work limit before rank
reduction.

The resulting finite-algebra status is \emph{formula verified} and
\emph{AI reviewed}. The residual-padded Schur gate remains operational
floating-point evidence, not a proved perturbation bound. No model review is
human-independent verification or an asymptotic theorem. Cluster independence,
realized-data scale, applicability, and named-human review remain open, so
every clustered result retains
\path{TALO_THEORY_APPLICABILITY_UNVERIFIED}.

\end{document}
